%% kapa: exam=1 round=34 year=2023 period=1 form=A subject=real_estate subject_ko=부동산학원론 seq=6 number=86
%% source: Q-Net 감정평가사 기출문제 (cst00309) artlSeq=5211789 file=raw/1008/5211789/2023년도 제34회 감정평가사 1차시험 1교시.pdf
%% conversion: kordoc 4.24.51 -> tools/gen1.py (auto); review=proofread-sonnet; flags=none
\begin{question}{86}
감정평가사 A는 표준지공시지가의 조사·평가를 의뢰받고 실지조사를 통해 표준지에 대해 다음과 같이 확인하였다. 표준지조사·평가보고서상 토지특성 기재방법의 연결이 옳은 것은?
\begin{framedtext}
ㄱ. 토지이용상황: 주변의 토지이용상황이 ‘전’으로서 돈사와 우사로 이용되고 있음 \par
ㄴ. 도로접면: 폭 10미터의 도로와 한면이 접하면서 자동차 통행이 불가능한 폭 2미터의 도로에 다른 한면이 접함 \par
\end{framedtext}
\begin{choices}
  \item ㄱ: 전기타, \quad ㄴ: 중로한면
  \item ㄱ: 전기타, \quad ㄴ: 소로한면
  \item ㄱ: 전축사, \quad ㄴ: 소로각지
  \item ㄱ: 전축사, \quad ㄴ: 소로한면
  \item ㄱ: 목장용지, \quad ㄴ: 소로한면
\end{choices}
\end{question}
\answer{4}
